% This is text associated with the open textbook "Open Electromagnetics"
% See immediately below for author, date
% This work is licensed under the Creative Commons Attribution-ShareAlike 4.0 International License. (CC BY-SA 4.0) 
% To view a copy of the license, visit http://creativecommons.org/licenses/by-sa/4.0/

%ORB TITLE Lossless and Low-Loss Transmission Lines
%ORB AUTHOR 0        # S.W. Ellingson, Virginia Tech (ellingson@vt.edu)
%ORB DATE 20180802   # yyyymmdd
%ORB ABSTRACT Simpler forms of alpha, beta, Z0 etc when line is lossless 

\label{m0083_Lossless_Transmission_Lines} % <-- Must match name of this file and module directory name.  Don't mess with this!

Quite often the loss in a transmission line is small enough that it may be neglected.  
In this case, several aspects of transmission line theory may be simplified.  
In this section, we present these simplifications.

First, recall that ``loss'' refers to the reduction of magnitude as a wave propagates through space. 
In the lumped-element equivalent circuit model 
\index{transmission line!lumped-element model}
(Section~\ref{m0029_Lumped-Element_Model}), the parameters $R'$ and $G'$ of the represent physical mechanisms associated with loss.
Specifically, $R'$ represents the resistance of conductors, whereas $G'$ represents the undesirable current induced between conductors through the spacing material.
Also recall that the propagation constant $\gamma$ is, in general, given by
%
\begin{equation}
\gamma \triangleq \sqrt{\left(R'+j\omega L'\right)\left(G'+j\omega C'\right)}
\end{equation} 

With this in mind, we now define ``low loss''
\index{transmission line!low loss}
as meeting the conditions:
%
\begin{equation}
R'\ll\omega L'
\label{m0083_eLLR}
\end{equation}
%
\begin{equation}
G'\ll\omega C'
\label{m0083_eLLG}
\end{equation}
%
When these conditions are met, the propagation constant simplifies as follows:
%
\begin{flalign}
\gamma &\approx \sqrt{\left(j\omega L'\right)\left(j\omega C'\right)} \nonumber \\
       &= \sqrt{-\omega^2 L'C'} \nonumber \\
       &= j\omega\sqrt{L'C'} 
\end{flalign} 
%
and subsequently
%
\begin{flalign}
\alpha \triangleq \mbox{Re}\left\{\gamma\right\} & \approx 0                     & \mbox{(low-loss approx.)} \\
\beta  \triangleq \mbox{Im}\left\{\gamma\right\} & \approx \omega\sqrt{L'C'}     & \mbox{(low-loss approx.)} \\
v_p    = \omega/\beta                            & \approx \frac{1}{\sqrt{L'C'}} & \mbox{(low-loss approx.)}
\end{flalign}
%
Similarly:
\begin{equation}
Z_0 = \sqrt{\frac{R' + j\omega L'}{G' + j\omega C'}} \approx \sqrt{\frac{L'}{C'}} ~~~ \mbox{(low-loss approx.)}
\end{equation} 

Of course if the line is strictly lossless (i.e., $R'=G'=0$) then these are not approximations, but rather the exact expressions.

In practice, these approximations are quite commonly used, since practical transmission lines typically meet the conditions expressed in Inequalities~\ref{m0083_eLLR} and \ref{m0083_eLLG} and the resulting expressions are much simpler.  We further observe that $Z_0$ and $v_p$ are approximately independent of frequency when these conditions hold.

However, also note that ``low loss'' does not mean ``no loss,'' and it is common to apply these expressions even when $R'$ and/or $G'$ is large enough to yield significant loss.  For example, a coaxial cable used to connect an antenna on a tower to a radio near the ground typically has loss that is important to consider in the analysis and design process, but nevertheless satisfies Equations~\ref{m0083_eLLR} and \ref{m0083_eLLG}.  In this case, the low-loss expression for $\beta$ is used, but $\alpha$ might not be approximated as zero.
 



